\uuid{aIgV}
\titre{Tangente}
\auteur{bodin ; hanani ; mzari}
\organisation{exo7}
\licence{CC BY-NC-SA 4.0}
\datecreate{2020-01-20}
\niveau{L1}
\module{Analyse}
\chapitre{Dérivabilité des fonctions réelles}
\sousChapitre{Autre}
\difficulte{2}
\theme{Dérivées, tangente}
\exo7qcmid{lille-1.derivees.018}

\contenu{
\begin{question}
\qclassification{124.00}
\qtags{motcle=tangente}
Soit $\displaystyle f(x)=\frac{x+3}{x+2}$. Quel est l'ensemble $S$ des points $x_0$ où la tangente à $\mathcal{C}_f$ est perpendiculaire à la droite d'équation $y=x$ ?
\begin{answers}
    \bad{$S=\{-2\}$}
    \bad{$S=\{-3\}$}
    \good{$S=\{-1,-3\}$}
    \bad{$S=\varnothing$}
\end{answers}
\begin{explanations}
La pente de la droite $y=x$ est $1$, donc la tangente à $\mathcal{C}_f$ en $x_0$ est perpendiculaire à cette droite si, et seulement si, $f'(x_0)=-1$. C'est-à-dire $x_0=-1$ ou $x_0=-3$.
\end{explanations}
\end{question}
}
